\documentclass[../../main.tex]{subfiles}
\begin{document}
% SOURCE p=231; printed=216; object=ch07-regularity-title; region=middle
\section{正则性与逼近定理}\label{section:ch07-regularity-approximation}\label{section:ch07-lp-regular-borel}
本节更详细地研究 Radon 测度的性质。

% SOURCE p=231; printed=216; object=ch07-sigma-finite-inner; number=7.5; region=middle
\begin{proposition}\label{proposition:ch07-sigma-finite-inner}
每个 Radon 测度在其所有 $\sigma$ 有限的 Borel 集上都内正则。
\end{proposition}
% SOURCE p=231; printed=216; object=ch07-sigma-finite-inner-proof; region=middle
\begin{proof}
设 $\mu$ 为 Radon 测度，$E$ 为 $\sigma$ 有限集。若 $\mu(E)<\infty$，给定 $\varepsilon>0$，可取开集 $U\supseteq E$ 使 $\mu(U)<\mu(E)+\varepsilon$，再取紧集 $F\subseteq U$ 使 $\mu(F)>\mu(U)-\varepsilon$。因为 $\mu(U\setminus E)<\varepsilon$，还可选开集 $V\supseteq U\setminus E$ 使 $\mu(V)<\varepsilon$。令 $K=F\setminus V$，则 $K$ 紧、$K\subseteq E$，并且
\begin{align*}
\mu(K)=\mu(F)-\mu(F\cap V)>\mu(E)-\varepsilon-\mu(V)>\mu(E)-2\varepsilon.
\end{align*}
所以 $\mu$ 在 $E$ 上内正则。若 $\mu(E)=\infty$，把 $E$ 写成有限测度 Borel 集 $E_j$ 的递增并，则 $\mu(E_j)\to\infty$。任给 $N\in\mathbb N$，取 $j$ 使 $\mu(E_j)>N$，由有限测度情形可取紧集 $K\subseteq E_j$ 使 $\mu(K)>N$。故 $\mu$ 在 $E$ 上也内正则。
\end{proof}
% SUPPLEMENT: id=v14-ch07-sigma-finite-inner-details; parent=ch07-sigma-finite-inner; kind=proof-detail+motivation+hypotheses; origin=independent-derivation
\begin{remark}
\textbf{两次外逼近的作用（重构新增）。}第一次把 $E$ 放进有限测度开集 $U$，才能用定义中仅对开集已知的内正则性取 $F$；第二次用开集 $V$ 除去 $F$ 中可能落在 $E$ 外的点。$F\setminus V$ 是紧集的闭子集，因而仍紧；若其中一点不在 $E$，它属于 $U\setminus E\subseteq V$，矛盾。选 $V$ 时，先利用严格不等式 $\mu(U\setminus E)<\varepsilon$，再在正的余量内作外逼近。上式只减有限量，没有 $\infty-\infty$。无限情形中，若 $E=\bigcup_j H_j$ 且各 $\mu(H_j)<\infty$，取 $E_j=\bigcup_{k\leqslant j}H_k$ 即得所需递增列。取紧集时使误差小于 $\mu(E_j)-N$，便得到严格的 $\mu(K)>N$。

动机是把开集上的内正则性传到可测集。Borel 条件保证各差集可测；Radon 的外正则性与开集内正则性分别用于上述步骤，紧集有限性保证紧近似本身具有有限质量。$\sigma$ 有限性使无限质量也能由可数个有限块恢复；它不能从一般 LCH 假设自动推出。\cref{exercise:ch07-uncountable-lines}将展示不满足这项条件时内正则性确实可能失败。
\end{remark}
% END SUPPLEMENT: id=v14-ch07-sigma-finite-inner-details

% SOURCE p=231; printed=216; object=ch07-sigma-finite-regular; number=7.6; region=middle-lower
\begin{corollary}\label{corollary:ch07-sigma-finite-regular}
每个 $\sigma$ 有限 Radon 测度都是正则测度。若 $X$ 为 $\sigma$ 紧空间，则 $X$ 上的每个 Radon 测度都正则。
\end{corollary}
% SUPPLEMENT: id=v14-ch07-sigma-finite-regular-proof; parent=ch07-sigma-finite-regular; kind=omitted-proof+motivation+hypotheses; origin=independent-derivation
\begin{proof}
\textbf{补全证明（重构新增）。}若 $X=\bigcup_j E_j$ 且 $\mu(E_j)<\infty$，每个 Borel 集 $E$ 都是 $E\cap E_j$ 的可数并，所以它 $\sigma$ 有限。由\cref{proposition:ch07-sigma-finite-inner}，$\mu$ 在每个 Borel 集上内正则；外正则已在 Radon 定义中，故 $\mu$ 正则。若 $X=\bigcup_j K_j$ 且各 $K_j$ 紧，则 $\mu(K_j)<\infty$，所以 $\mu$ 为 $\sigma$ 有限，应用第一部分即可。
\end{proof}
\begin{remark}
\textbf{动机与条件（重构新增）。}$\sigma$ 紧性是拓扑条件，$\sigma$ 有限性是关于测度的条件；Radon 的紧集有限性把前者转换成后者。结论不要求总测度有限，也不要求空间有可数基。
\end{remark}
% END SUPPLEMENT: id=v14-ch07-sigma-finite-regular-proof
非正则的 Radon 测度的例子见\cref{exercise:ch07-uncountable-lines}。

% SOURCE p=231; printed=216; object=ch07-open-closed-approximation; number=7.7; region=lower
\begin{proposition}\label{proposition:ch07-open-closed-approximation}
设 $\mu$ 为 $X$ 上的 $\sigma$ 有限 Radon 测度，$E\in\mathscr B_X$。
\begin{enumerate}[label=\alph*.,ref=(\alph*)]
\item\label{item:ch07-open-closed-epsilon} 对每个 $\varepsilon>0$，存在开集 $U$ 和闭集 $F$，使 $F\subseteq E\subseteq U$ 且 $\mu(U\setminus F)<\varepsilon$。
\item\label{item:ch07-fsigma-gdelta-null} 存在 $F_\sigma$ 集 $A$ 和 $G_\delta$ 集 $B$，使 $A\subseteq E\subseteq B$ 且 $\mu(B\setminus A)=0$。
\end{enumerate}
\end{proposition}
% SOURCE p=231; printed=216; object=ch07-open-closed-approximation-proof; region=bottom
\begin{proof}
写 $E=\bigcup_{j=1}^\infty E_j$，其中 $E_j$ 两两不交且测度有限。为每个 $j$ 选开集 $U_j\supseteq E_j$ 使 $\mu(U_j)<\mu(E_j)+\varepsilon2^{-j-1}$，令 $U=\bigcup_jU_j$。则 $U$ 开、$U\supseteq E$，并有
\begin{align*}
\mu(U\setminus E)\leqslant\sum_{j=1}^\infty\mu(U_j\setminus E_j)<\varepsilon/2.
\end{align*}
对 $E^c$ 作同样论证，得到开集 $V\supseteq E^c$ 且 $\mu(V\setminus E^c)<\varepsilon/2$。令 $F=V^c$，则 $F$ 闭、$F\subseteq E$，且
\begin{align*}
\mu(U\setminus F)=\mu(U\setminus E)+\mu(E\setminus F)
=\mu(U\setminus E)+\mu(V\setminus E^c)<\varepsilon.
\end{align*}
这证明了 (a)；(b) 容易推出，原书将细节留给读者。
\end{proof}
% SUPPLEMENT: id=v14-ch07-open-closed-details; parent=ch07-open-closed-approximation; kind=omitted-proof+proof-detail+motivation+hypotheses; origin=independent-derivation
\begin{remark}
\textbf{(b) 的完整证明及条件用途（重构新增）。}对每个 $n$，在 (a) 中取误差 $2^{-n}$，得到闭集 $F_n\subseteq E$ 及开集 $U_n\supseteq E$。令 $A=\bigcup_nF_n$、$B=\bigcap_nU_n$，则它们分别为 $F_\sigma$、$G_\delta$ 集，并满足所需包含关系。对每个 $n$，$B\setminus A\subseteq U_n\setminus F_n$，故 $0\leqslant\mu(B\setminus A)<2^{-n}$；令 $n\to\infty$ 得其测度为零。无需假设两列本身嵌套。

动机是在集合可能测度无穷时仍直接控制“夹层”的测度。(a) 的闭集 $F$ 未必紧；在无穷质量集合内通常不能以一个紧集留下任意小误差。全空间 $\sigma$ 有限性同时用于分割 $E$ 和 $E^c$，不是只要求 $E$ 有限。可先把一个可数有限测度覆盖不交化，再与目标集合相交，即得到证明中的 $E_j$。正则逼近与可求和误差相配合，不曾对无限测度作减法。
\end{remark}
% END SUPPLEMENT: id=v14-ch07-open-closed-details

% SOURCE p=232; printed=217; object=ch07-automatic-regularity; number=7.8; region=top
\begin{theorem}\label{theorem:ch07-automatic-regularity}
设 $X$ 为 LCH 空间，且每个开集都 $\sigma$ 紧（例如 $X$ 第二可数时如此）。那么 $X$ 上每个在紧集上有限的 Borel 测度都正则，因而都是 Radon 测度。
\end{theorem}
% SOURCE p=232; printed=217; object=ch07-automatic-regularity-proof; region=upper
\begin{proof}
设 $\mu$ 为这样的测度，则 $C_c(X)\subseteq L^1(\mu)$，所以 $I(f)=\int f\,\mathrm d\mu$ 为正线性泛函。由\cref{theorem:ch07-riesz-positive}，它对应一个 Radon 测度 $\nu$。若 $U$ 开，把它写为紧集的并 $U=\bigcup_{j=1}^\infty K_j$。先选 $f_1\prec U$ 且 $f_1=1$ 于 $K_1$；递归地选 $f_n\prec U$，使它在
\begin{align*}
\bigcup_{j=1}^nK_j\ \cup\ \bigcup_{j=1}^{n-1}\operatorname{supp}(f_j)
\end{align*}
上等于一。则 $f_n$ 逐点递增到 $\chi_U$，故由单调收敛定理，
\begin{align*}
\mu(U)=\lim_n\int f_n\,\mathrm d\mu
=\lim_n\int f_n\,\mathrm d\nu=\nu(U).
\end{align*}
给定 Borel 集 $E$ 及 $\varepsilon>0$，由\cref{proposition:ch07-open-closed-approximation}，存在开集 $V\supseteq E$、闭集 $F\subseteq E$ 使 $\nu(V\setminus F)<\varepsilon$。$V\setminus F$ 开，故 $\mu(V\setminus F)=\nu(V\setminus F)<\varepsilon$。于是 $\mu(V)\leqslant\mu(E)+\varepsilon$，证明外正则性；并且 $\mu(F)\geqslant\mu(E)-\varepsilon$。$F$ 为 $\sigma$ 紧集，因为 $X$ 为 $\sigma$ 紧空间；取递增紧集 $L_j\subseteq F$ 并满 $F$，则 $\mu(L_j)\to\mu(F)$，证明内正则性。因此 $\mu$ 正则，并由表示定理的唯一性有 $\mu=\nu$。
\end{proof}
% SUPPLEMENT: id=v14-ch07-automatic-regularity-details; parent=ch07-automatic-regularity; kind=proof-detail+motivation+hypotheses; origin=independent-derivation
\begin{remark}
\textbf{可数性与无限值的核验（重构新增）。}每一步要求函数等于一的集合都是 $U$ 内的有限个紧集之并，故局部 Urysohn 引理适用。前一个函数的支集上新函数等于一，支集外前者为零，所以确有 $f_n\geqslant f_{n-1}$。$U$ 的每一点终落入某个 $K_j$，于是极限是 $\chi_U$。

应用开闭逼近命题前须核对 $\nu$ 的 $\sigma$ 有限性：$X$ 本身是开集，故 $\sigma$ 紧；$\nu$ 在紧集上有限，因而 $\sigma$ 有限。若 $X=\bigcup_j H_j$ 且 $H_j$ 紧，则 $L_j=F\cap\bigcup_{k\leqslant j}H_k$ 是递增紧集并满 $F$。有限 $\mu(E)$ 时上述误差直接给出内外正则性；若 $\mu(E)=\infty$，外正则自动成立，且由 $\mu(E\setminus F)<\varepsilon$ 可知 $\mu(F)=\infty$，从而 $\mu(L_j)\to\infty$。这是无穷情形的严格解释，不使用 $\infty-\infty$。

动机是利用连续测试把未知测度与已知 Radon 测度比较。紧集有限使测试可积；\emph{每个开集} $\sigma$ 紧才允许用序列恢复其测度，仅知全空间 $\sigma$ 紧不能完成这一步。第二可数 LCH 空间满足此条件：在任意开集内用相对紧开邻域覆盖，再利用该开子空间的第二可数性选可数子覆盖；其闭包都是开集内的紧集，这也是\subcref{exercise:ch07-second-countable-baire}{item:ch07-second-countable-baire-borel}中已证明的论证。
\end{remark}
% END SUPPLEMENT: id=v14-ch07-automatic-regularity-details

% SOURCE p=232; printed=217; object=ch07-nonradon-examples-transition; region=middle
非 Radon 测度的例子见\cref{exercise:ch07-zero-weight-failure,exercise:ch07-semifinite-radon,exercise:ch07-ordinal-measure}。特别地，最后一题给出紧 Hausdorff 空间上的一个有限但非 Radon 的 Borel 测度。现在转向可测函数的逼近定理。

% SOURCE p=232; printed=217; object=ch07-cc-lp-density; number=7.9; region=middle-lower
\begin{proposition}\label{proposition:ch07-cc-lp-density}
若 $\mu$ 为 $X$ 上的 Radon 测度，则当 $1\leqslant p<\infty$ 时，$C_c(X)$ 在 $L^p(\mu)$ 中稠密。
\end{proposition}
% SOURCE p=232; printed=217; object=ch07-cc-lp-density-proof; region=middle-lower
\begin{proof}
由\cref{proposition:ch06-simple-density}，$L^p$ 简单函数在 $L^p$ 中稠密，因此只须对任意有限测度 Borel 集 $E$ 以 $C_c(X)$ 中函数按 $L^p$ 范数逼近 $\chi_E$。给定 $\varepsilon>0$，由\cref{proposition:ch07-sigma-finite-inner}及外正则性，取紧集 $K\subseteq E$ 及开集 $U\supseteq E$ 使 $\mu(U\setminus K)<\varepsilon$。由 Urysohn 引理取 $f\in C_c(X)$ 满足 $\chi_K\leqslant f\leqslant\chi_U$，则
\begin{align*}
\|\chi_E-f\|_p\leqslant\mu(U\setminus K)^{1/p}<\varepsilon^{1/p}.
\end{align*}
这就完成证明。
\end{proof}
% SUPPLEMENT: id=v14-ch07-cc-lp-density-details; parent=ch07-cc-lp-density; kind=proof-detail+motivation+hypotheses; origin=independent-derivation
\begin{remark}
\textbf{从集合逼近到范数逼近（重构新增）。}分别把内外误差选为小于 $\varepsilon/2$，即得 $\mu(U\setminus K)<\varepsilon$。在 $K$ 内两个函数都等于一，在 $U^c$ 上都为零，其余点之差的模不超过一，所以 $|\chi_E-f|^p\leqslant\chi_{U\setminus K}$。对有限和 $s=\sum_{j=1}^m c_j\chi_{E_j}$，逐项选近似 $f_j$，三角不等式给出 $\|s-\sum_j c_jf_j\|_p\leqslant\sum_j|c_j|\|\chi_{E_j}-f_j\|_p$；分配足够小的误差便完成向一般简单函数及一般 $L^p$ 函数的过渡。

Radon 条件为有限测度集提供紧内逼近与开外逼近，LCH 条件提供两者之间的连续截断。只需各目标简单函数的支撑测度有限，不需全空间 $\sigma$ 有限。有限 $p$ 使小测度误差转化为小范数误差；$p=\infty$ 时不成立，例如 Lebesgue 测度实线上 $\chi_{[0,1]}$ 与任意连续函数的本质上确界距离至少为 $1/2$：若距离小于 $1/2$，在零点两侧分别由几乎处处的严格界及连续性，强迫函数在零点同时小于和大于 $1/2$（先选一个共同的误差界 $c<1/2$ 即可）。若采用测度的完备化，所有可积等价类均有 Borel 代表，上述稠密结论保持成立。
\end{remark}
% END SUPPLEMENT: id=v14-ch07-cc-lp-density-details

% SOURCE p=232,233; printed=217-218; object=ch07-lusin; number=7.10; region=bottom-top
\begin{theorem}[Lusin 定理]\label{theorem:ch07-lusin}
设 $\mu$ 为 $X$ 上的 Radon 测度，$f:X\to\mathbb C$ 可测，且在某个有限测度集外为零。则对任意 $\varepsilon>0$，存在 $\phi\in C_c(X)$，使 $\phi=f$ 除了一个测度小于 $\varepsilon$ 的集合。若 $f$ 有界，还可要求 $\|\phi\|_u\leqslant\|f\|_u$。
\end{theorem}
% SOURCE p=232,233; printed=217-218; object=ch07-lusin-proof; region=bottom-upper
\begin{proof}
令 $E=\{x:f(x)\ne0\}$，先设 $f$ 有界，则 $f\in L^1(\mu)$。由\cref{proposition:ch07-cc-lp-density}，存在 $g_n\in C_c(X)$ 在 $L^1$ 中收敛到 $f$，再由\cref{corollary:ch02-lone-ae-subsequence}取子列（仍记为 $g_n$）几乎处处收敛到 $f$。由 Egoroff 定理，存在 $A\subseteq E$ 使 $\mu(E\setminus A)<\varepsilon/3$，并且 $g_n\to f$ 在 $A$ 上一致收敛。再取紧集 $B\subseteq A$ 及开集 $U\supseteq E$ 使
\begin{align*}
\mu(A\setminus B)<\varepsilon/3,\qquad \mu(U\setminus E)<\varepsilon/3.
\end{align*}
一致收敛说明 $f|_B$ 连续；由\cref{theorem:ch04-local-tietze}，存在 $h\in C_c(X)$，在 $B$ 上等于 $f$，且 $\operatorname{supp}(h)\subseteq U$。于是 $\{f\ne h\}\subseteq U\setminus B$，后者测度小于 $\varepsilon$。

令 $M=\|f\|_u$，定义 $\beta:\mathbb C\to\mathbb C$：当 $|z|\leqslant M$ 时 $\beta(z)=z$，当 $|z|>M$ 时 $\beta(z)=M\operatorname{sgn}z$。令 $\phi=\beta\circ h$。$\beta$ 连续且 $\beta(0)=0$，故 $\phi\in C_c(X)$；同时 $\|\phi\|_u\leqslant M$，并且在 $h=f$ 处仍有 $\phi=f$。

若 $f$ 无界，令 $A_n=\{x:0<|f(x)|\leqslant n\}$。则 $A_n\uparrow E$，故充分大 $n$ 满足 $\mu(E\setminus A_n)<\varepsilon/2$。对有界函数 $f\chi_{A_n}$ 应用已证结论，得到 $\phi\in C_c(X)$ 在一个测度小于 $\varepsilon/2$ 的集合外等于该函数。加上 $E\setminus A_n$，便知 $\phi=f$ 的例外集测度小于 $\varepsilon$。
\end{proof}
% SUPPLEMENT: id=v14-ch07-lusin-details; parent=ch07-lusin; kind=proof-detail+motivation+hypotheses; origin=independent-derivation
\begin{remark}
\textbf{逐条件核验与动机（重构新增）。}这里 $f$ 取有限复值，所以 $A_n\uparrow E$；若允许无穷值，截断后的相等结论必须另作处理。$\mu(E)<\infty$ 同时保证有界时的可积性、Egoroff 定理的适用以及 $\mu(E\setminus A_n)\to0$。由 $B\subseteq A\subseteq E\subseteq U$，$U\setminus B$ 分成 $U\setminus E$、$E\setminus A$、$A\setminus B$ 三个互不相交部分，因此误差预算严格小于 $\varepsilon$。在 $B$ 上，一致极限连续；紧性与 LCH 条件才允许把这个局部连续函数延拓为支集位于 $U$ 内的全局连续函数。复值延拓也可对实、虚部分分别应用局部 Tietze 定理再合并。径向截断保持闭圆盘内每一点，故不破坏相等点；它保持零值，故不扩大支集。$M=0$ 时 $\beta\equiv0$，无除零问题。

动机是在大部分点上实现\emph{精确相等}，而不只是范数意义下接近。外侧 $U$ 控制延拓可能新产生的非零点，内侧 $B$ 保留函数的原值，两者缺一不可。本章按 Borel 可测理解可测；若对完备化工作，可先取 Borel 代表，原函数与代表相异的零测度集并入例外集。这样仍可得到严格小于 $\varepsilon$ 的误差。
\end{remark}
% END SUPPLEMENT: id=v14-ch07-lusin-details

\subfile{section-02-semicontinuity.tex}
\end{document}
